\section{Tốc độ phản ứng và chất xúc tác}

Phần tốc độ phản ứng trong chương trình lớp 8 chủ yếu yêu cầu nhận biết phản ứng nhanh hay chậm, nêu được một số yếu tố làm thay đổi tốc độ và hiểu vai trò của chất xúc tác. Vì vậy, trong chuyên đề này phần tốc độ được giữ ở mức cốt lõi; các công thức động học nâng cao không phải trọng tâm.

\subsection{Khái niệm và cách so sánh}

\textbf{Tốc độ phản ứng} cho biết phản ứng xảy ra nhanh hay chậm. Có thể so sánh định tính bằng cách:
\begin{itemize}
\item trong cùng một thời gian, phản ứng nào tạo nhiều sản phẩm hơn thì nhanh hơn;
\item để tạo cùng một lượng sản phẩm, phản ứng nào cần ít thời gian hơn thì nhanh hơn.
\end{itemize}

\begin{ghinho}{Các yếu tố thường gặp}
\begin{itemize}
\item \textbf{Nồng độ:} nồng độ chất phản ứng lớn hơn thường làm phản ứng nhanh hơn.
\item \textbf{Nhiệt độ:} tăng nhiệt độ thường làm phản ứng nhanh hơn.
\item \textbf{Diện tích bề mặt:} nghiền nhỏ chất rắn làm tăng diện tích tiếp xúc nên phản ứng thường nhanh hơn.
\item \textbf{Chất xúc tác:} làm phản ứng diễn ra nhanh hơn nhưng không bị tiêu hao theo phương trình tổng quát.
\end{itemize}
\end{ghinho}

\begin{vidu}
Cho hai mẫu \ce{CaCO3} có cùng khối lượng, một mẫu dạng cục và một mẫu nghiền nhỏ, lần lượt tác dụng với cùng thể tích dung dịch \ce{HCl} cùng nồng độ và cùng nhiệt độ. Mẫu nào phản ứng nhanh hơn?
\end{vidu}
\begin{loigiai}
Mẫu nghiền nhỏ có tổng diện tích bề mặt tiếp xúc với dung dịch acid lớn hơn. Vì vậy, trong cùng thời gian có nhiều vị trí phản ứng hơn.

\textbf{Vậy,} mẫu \ce{CaCO3} nghiền nhỏ phản ứng nhanh hơn mẫu dạng cục.
\end{loigiai}

\begin{vidu}
Trong phản ứng phân hủy hydrogen peroxide:
\[
\ce{2H2O2 ->[MnO2] 2H2O + O2},
\]
khi thêm một lượng nhỏ \ce{MnO2}, khí oxygen thoát ra nhanh hơn. Nêu vai trò của \ce{MnO2}.
\end{vidu}
\begin{loigiai}
\ce{MnO2} làm phản ứng xảy ra nhanh hơn và không bị tiêu hao theo phương trình tổng quát.

\textbf{Vậy,} \ce{MnO2} là chất xúc tác của phản ứng phân hủy hydrogen peroxide.
\end{loigiai}

\subsection{Đọc và xử lí số liệu tốc độ phản ứng}

Khi đề cho bảng thể tích khí theo thời gian, có thể so sánh tốc độ bằng lượng khí tạo ra trong cùng một khoảng thời gian. Nếu đề yêu cầu tính toán, có thể dùng:
\[
\frac{\Delta V}{\Delta t}.
\]
Đây chỉ là công cụ đọc dữ liệu thực nghiệm, không phải trọng tâm tính toán của chuyên đề.

\begin{vidu}
Hai thí nghiệm A và B đều tạo cùng loại khí. Sau $30$ s, A thu được $45$ mL khí, B thu được $70$ mL khí. Sau thời gian đủ dài, cả hai đều thu được $100$ mL khí. Hãy so sánh tốc độ ban đầu và lượng sản phẩm cuối cùng.
\end{vidu}
\begin{loigiai}
Trong cùng $30$ s, B tạo nhiều khí hơn nên B phản ứng nhanh hơn ở giai đoạn đầu. Tuy nhiên, khi phản ứng hoàn tất, cả hai đều thu được $100$ mL khí.

\textbf{Vậy,} B có tốc độ ban đầu lớn hơn, nhưng lượng khí cuối cùng của hai thí nghiệm bằng nhau.
\end{loigiai}
